% Generated draft: requires mathematical and source review.
\documentclass[11pt]{article}
\usepackage[margin=25mm]{geometry}
\usepackage{fontspec}
\setmainfont{Latin Modern Roman}
\usepackage{amsmath,amssymb,mathtools,unicode-math}
\setmathfont{Latin Modern Math}
\usepackage{enumitem,needspace,xurl,hyperref}
\setlength{\parindent}{0pt}
\setlength{\parskip}{5pt}
\setlength{\emergencystretch}{3em}
\newcommand{\N}{\mathbb{N}}
\newcommand{\Z}{\mathbb{Z}}
\newcommand{\Q}{\mathbb{Q}}
\newcommand{\R}{\mathbb{R}}
\newcommand{\C}{\mathbb{C}}
\newcommand{\F}{\mathbb{F}}
\newcommand{\A}{\mathbb{A}}
\newcommand{\PP}{\mathbb{P}}
\newcommand{\E}{\mathbb{E}}
\newcommand{\eps}{\varepsilon}
\newcommand{\dd}{\,\mathrm{d}}
\begin{document}
\section*{UnsolvedMath ID: 30005699}
\textbf{Source problem code:} catalogue entry
\subsection{Definitions and mathematical statement}
A finitely generated free group is a group isomorphic to the free group \(F_n\) on a finite basis of size \(n\). The original source describes the groups under consideration as \(\{\text{finitely generated free}\}\)-by-cyclic groups. In extension notation, this means groups \(G\) admitting an exact sequence
\[
1\longrightarrow F_n\longrightarrow G\longrightarrow C\longrightarrow 1,
\]
where \(F_n\) is finitely generated and free and \(C\) is cyclic. The supplied excerpt does not specify whether ``cyclic'' means infinite cyclic or also permits finite cyclic quotients; the intended range of \(C\) must therefore be checked before this becomes a fully specified quantified statement.

\textbf{Conjecture 6 (as supplied).} Does every \(\{\text{finitely generated free}\}\)-by-cyclic group \(G\) belong to
\[
\operatorname{TAP}_{\infty}(R)?
\]
The supplied record does not define \(\operatorname{TAP}_{\infty}(R)\), specify the nature or hypotheses of the parameter \(R\), or state its quantifier. Membership in this class is consequently retained as source notation, without an invented definition. The question is a statement schema pending recovery of those definitions and conventions. In particular, the original wording requires the free normal subgroup to be finitely generated; this requirement must not be replaced merely by finite generation of \(G\).
\subsection{Short English statement}
The supplied Conjecture 6 asks whether every group with a finitely generated free normal subgroup and cyclic quotient lies in \(\operatorname{TAP}_{\infty}(R)\). A complete formulation requires the source definition of this class, the assumptions and quantifier on \(R\), and the convention for cyclic quotients. The registry labels the problem open; this is supplied curation, not an independently verified status.
\subsection{Sources}
\begin{itemize}
\item Supplied local UnsolvedMath registry record, permanent ID \(30005699\), title ``TAP Infinity for Free-by-Cyclic Groups,'' source problem code \texttt{catalogue entry}, in the Oberwolfach Reports Open Problems collection. No individual proposer is identified in the supplied record.
\item Supplied source citation: \emph{Mini-Workshop: Homological aspects of TDLC-groups} (2024), Oberwolfach Reports, pp.\ 3103--3138, DOI \texttt{10.4171/owr/2023/55}: \url{https://doi.org/10.4171/owr/2023/55}. The original excerpt identifies Conjecture 6 on p.\ 3127. The registry gives proposed year 2023; the supplied citation gives 2024.
\item Supplied literature source, ``Oberwolfach report'': \url{https://ems.press/content/serial-article-files/48166}.
\item Supplied literature source, ``Hughes--Kudlinska'': \url{https://doi.org/10.1112/plms.70059}.
\item Supplied literature source, ``Profinite rigidity of fibring'': \url{https://arxiv.org/abs/2206.11347}.
\item Literature assessment dated exactly \texttt{2026-08-22}: supplied dated curation, not independently verified. It labels the problem open and states that known TAP results cover LERF and products of limit groups, while the 2025 free-by-cyclic work does not prove TAP infinity for every finitely generated free-by-cyclic group. The background explicitly requests independent verification of time-sensitive status and cited claims.
\item The registry labels the statement \texttt{corrected\_verified} and reports that it was checked against the official Oberwolfach report. This is a supplied verification claim; no source consultation or independent mathematical verification is asserted here.
\end{itemize}
\end{document}
